% !PRINTED 208
\weberpage{208}
%<<CONT
is replaced by another, by replacing $\nu'$ in (5) by another symbol, $m$, and then setting $nu = 2\nu'\omega$. If we restrict ourselves to the case $(g_1, g_2) = (0, 1)$ and further multiply (5) by $\vartheta_{00} : \vartheta_{10}$, there follows:

\weq{12}{
\sum^{\nu} e^{\frac{8m\nu\pi i}{n}} x_{\nu,\nu'} = \\
(-1)^{\frac{n-1}{2}} n\, e^{\frac{-4\pi i m\nu'\omega}{n}} \frac{\vartheta_{00}\,\vartheta_{01}\,\vartheta'_{11}(0, n\omega)\,\vartheta_{11}[2(\nu'-m)\omega, n\omega]}{\vartheta_{10}\,\vartheta'_{11}\,\vartheta_{01}(2m\omega, n\omega)\,\vartheta_{01}(2\nu'\omega, n\omega)},
}
% Source verified: the summation sign carries only the upper index \nu.

where the left side vanishes \gesperrt{if and only if} $m \equiv \nu' \pmod{n}$.

\medskip

\begin{center}
\textbf{\S\,63. The \gesperrt{Galois Group} of the Division Equation.}
\end{center}

In the thirteenth chapter of the first volume it was shown how the algebraic nature of an equation finds its simplest expression in the \gesperrt{Galois group} of the equation.

Let the definition of the \gesperrt{Galois group} be briefly recalled here.

Once the domain of rationality $\Omega$ has been fixed, the \gesperrt{Galois group} of an equation $F(x) = 0$, of which it is assumed only that it has no repeated roots, is defined as the group $G$ of the permutations among the roots of this equation to which the twofold property belongs:

a) Every rational equation in $\Omega$ which subsists among the roots of $F(x)$ remains correct when the roots are subjected to any permutation of this group.

b) Every rational function in $\Omega$ of the roots of $F(x)$ which admits all the permutations of the group is contained in $\Omega$.

Now, as we seek to determine the group $G$ of the division equation, we assume as domain of rationality in the first place the \gesperrt{totality of all rational numbers and rational functions of $\varkappa^2$}.

By the addition and multiplication theorem (cf. \S\, 57, 61), if $f$ and $\varphi_m$ denote rational functions which are independent of $\mu, \mu', \nu, \nu'$, then % [source typo] The German source omits the comma between \mu' and \nu.

\weq{1}{x_{\mu+\nu,\,\mu'+\nu'} = f(x_{\mu,\mu'}, x_{\nu,\nu'}),}

\weq{2}{x_{m\mu,\,m\mu'} = \varphi_m(x_{\mu,\mu'}).}
